% ============================================================================ 0 · Cómo leer + criterios
\section*{Cómo leer este informe}
\addcontentsline{toc}{section}{Cómo leer este informe}

El informe sigue las siete secciones del documento de Ford. En cada una contamos qué entendimos, qué
encontramos en los datos y qué hicimos.

\begin{itemize}
  \item \tagtest{}: medido sobre los 111 autos del holdout, sorteados antes de mirar ninguna feature contra la
    etiqueta (103 con datos: 32 fallados y 71 sanos). Cada modelo se entrenó con todo el conjunto de desarrollo.
    \tagdev{}: validación cruzada agrupada por vehículo sobre los 446 autos de desarrollo. La demo, la explicabilidad
    y las auditorías son de dev.
  \item \textbf{Detección al X\,\%}: de los autos que fallan, cuántos reciben una alerta antes de la falla con un umbral
    fijado para que solo el X\,\% de los sanos tenga una falsa alarma. Una alerta son dos cortes seguidos sobre el umbral.
  \item \textbf{Azar}: el mismo puntaje permutado entre filas, conservando el largo del historial de cada auto.
  \item \textbf{Nada de accuracy}: con \textasciitilde6\,\% de filas positivas, “nadie falla” da 94\,\%. El PR-AUC por fila
    va al Anexo~\ref{anx:auditorias}, al lado del techo de cohorte.
\end{itemize}

\begin{table}[!htbp]
\centering
\caption{Los criterios de Ford (§5.1 y §5.2 del desafío) y dónde se responden.}
\label{tab:criterios}
\small
\begin{tabularx}{\textwidth}{@{}>{\raggedright\arraybackslash}p{3.4cm}>{\raggedright\arraybackslash}X>{\raggedright\arraybackslash}p{1.9cm}@{}}
\toprule
Criterio & Qué mostramos & Dónde \\
\midrule
Objetivo 5.1: anticipar con el mayor margen & \textasciitilde4.600 km (\textasciitilde100 días) de mediana en test, con
  500 km ciegos antes del evento; el puntaje sube hacia el evento & §\ref{sec:panel}, §\ref{sec:anticipacion} \\
Viabilidad técnica & Medido en un holdout congelado; GRU de \textasciitilde3.900 parámetros en CPU; demo funcionando;
  plan por etapas con piloto en sombra & §\ref{sec:viabilidad} \\
Costo & \textasciitilde USD 0,07–0,10 por auto y por año, sin sensores nuevos & §\ref{sec:costo} \\
Carácter innovador & Una métrica que separa \emph{qué auto} de \emph{cuándo}; red recurrente sobre la ventana; IA
  generativa con verificador en código; la perilla de falsas alarmas & §\ref{sec:innovacion} \\
Calidad del tratamiento de datos & Trampas de formato y de muestreo corregidas; atípicos con topes físicos; faltantes
  estructurales tratados como información; correcciones declaradas, holdout congelado y auditorías de leakage &
  §\ref{sec:datos}, Anexo~\ref{anx:auditorias} \\
Claridad para mostrar los datos & Curva de detección contra falsas alarmas con el azar dibujado; figuras con fuente;
  demo con tres pantallas y perilla & §\ref{sec:claridad} \\
Explicabilidad & En qué se apoya la GRU (permutación) y qué se le dice al conductor (comparación con la flota sana,
  filtrada por la física del DPF) & §\ref{sec:explicabilidad} \\
Pitch & La misma historia y los mismos números, en siete bloques y 29 minutos (\repo{docs/\allowbreak{}pitch/\allowbreak{}}) &
  §\ref{sec:claridad} \\
\bottomrule
\end{tabularx}
\end{table}
